\chapter{Einleitung}



In dieser Aufgabe baust du deinen eigenen digitalen Signalgenerator direkt auf 
dem FPGA. Kein fertiger Funktionsgenerator, kein BlackBox Chip – du entwickelst 
selbst die komplette Logik, die später echte Wellenformen erzeugt. Das heißt: 
Du steuerst Frequenz, Amplitude und Signalform live über Taster und siehst auf 
dem Oszilloskop sofort, was deine Schaltung macht. Genau so, wie man es sich 
als Entwickler eigentlich immer wünscht: Knopf drücken, Signal ändert sich, 
und du weißt, dass es \emph{deine} digitale Hardware ist, die das gerade tut.\\

Damit das klappt, kombinierst du mehrere richtig coole Bausteine: einen 
DDS basierten Sinusgenerator, sowie einen Rechteck und Sägezahn und Dreieckgenerator mit einen dynamischen Taktteiler. 
Über LEDs bekommst du jederzeit 
Feedback, welcher Parameter aktiv ist und welche Wellenform gerade ausgegeben 
wird. Wenn du alles sauber programmierst, kannst du per Knopfdruck zwischen 
Sinus, Dreieck, Rechteck und Sägezahn umschalten und die Parameter in Echtzeit 
verändern – ohne irgendeine Software, komplett in Hardware.\\

Kurz gesagt: Du baust dir einen kleinen, aber vollwertigen digitalen 
Funktionsgenerator – komplett selbst programmiert, komplett auf dem FPGA, 
und du siehst dein Ergebnis direkt auf dem Oszilloskop. Klingt gut. Dann los.